% Part: first-order-logic
% Chapter: sequent-calculus
% Section: proving-things-quant

\documentclass[../../../include/open-logic-section]{subfiles}

\begin{document}

\olfileid{fol}{seq}{prq}

\olsection{\usetoken{P}{derivation} met kwantoren (`voor alle' en `er bestaat')}

\begin{ex}
Geef een $\Log{LK}$-!!{derivation} van de sequent
$\lexists[x][\lnot !A(x)] \Sequent \lnot \lforall[x][!A(x)]$.

Bij kwantoren moeten we de eigenvariabelevoorwaarde in de gaten
houden. Die voorwaarde beperkt welke letter we in sommige regels mogen
invullen. Daardoor kan de volgorde van onze stappen ertoe doen. Pak de
regels waarvoor die voorwaarde geldt het liefst als eerste aan. In het
afgemaakte bewijs komen zulke regels dan lager te staan.
Daarnaast helpt het om vooruit te kijken naar de beginsequent waarop
we willen uitkomen. Hier moet dat iets als $!A(a) \Sequent !A(a)$ zijn.
Als we de kwantorregels van onder naar boven gebruiken, moeten we dus
dezelfde term kiezen. We noemen die term $a$ en gebruiken hem bij zowel
de $\lforall$- als de $\lexists$-regel. Met verschillende termen zouden
we iets als $!A(a) \Sequent !A(b)$ krijgen. Die sequent is niet
afleidbaar.

We beginnen zoals altijd met de sequent die onderaan moet staan:
\begin{prooftree}
\AxiomC{}
\UnaryInf$\lexists[x][\lnot !A(x)] \fCenter \lnot \lforall[x][!A(x)]$
\end{prooftree}
We kunnen nu \LeftR{\exists} of \RightR{\lnot} gebruiken. Voor
\LeftR{\exists} geldt de eigenvariabelevoorwaarde. Daarom handelen we
die regel meteen af en passen we hem als eerste toe.
\begin{prooftree}
\AxiomC{}
\UnaryInf$ \lnot !A(a) \fCenter \lnot \lforall[x][!A(x)]$
\RightLabel{\LeftR{\lexists}}
\UnaryInf$ \lexists[x][\lnot !A(x)] \fCenter \lnot \lforall[x][!A(x)]$
\end{prooftree}
Als we daarna \LeftR{\lnot} en \RightR{\lnot} van onder naar boven
toepassen, krijgen we
\begin{prooftree}
\AxiomC{}
\UnaryInf$\lforall[x][!A(x)] \fCenter !A(a)$
\RightLabel{\LeftR{\lnot}}
\UnaryInf$\lnot !A(a), \lforall[x][!A(x)] \fCenter $
\RightLabel{\LeftR{\Exchange}}
\UnaryInf$\lforall[x][!A(x)], \lnot !A(a) \fCenter $
\RightLabel{\RightR{\lnot}}
\UnaryInf$ \lnot !A(a) \fCenter \lnot \lforall[x] !A(x)$
\RightLabel{\LeftR{\lexists}}
\UnaryInf$ \lexists[x] \lnot !A(x) \fCenter \lnot \lforall[x] !A(x)$
\end{prooftree}
Nu kunnen we alleen nog \LeftR{\forall} toepassen. Voor deze regel geldt
de eigenvariabelebeperking niet. We willen uitkomen op een
beginsequent van de vorm $!A(a) \Sequent !A(a)$. Daarom vullen we bij
deze regel $a$ in als argument van $!A$.
\begin{prooftree}
\Axiom$!A(a) \fCenter !A(a)$
\RightLabel{\LeftR{\lforall}}
\UnaryInf$\lforall[x][!A(x)] \fCenter !A(a)$
\RightLabel{\LeftR{\lnot}}
\UnaryInf$\lnot !A(a), \lforall[x][!A(x)] \fCenter $
\RightLabel{\LeftR{\Exchange}}
\UnaryInf$\lforall[x][!A(x)], \lnot !A(a) \fCenter $
\RightLabel{\RightR{\lnot}}
\UnaryInf$ \lnot !A(a) \fCenter \lnot \lforall[x][!A(x)]$
\RightLabel{\LeftR{\lexists}}
\UnaryInf$ \lexists[x][ \lnot !A(x)] \fCenter \lnot \lforall[x][!A(x)]$
\end{prooftree}
Bij kwantoren moeten we aan het einde nog eens controleren of we de
eigenvariabelevoorwaarde nergens hebben geschonden. Die voorwaarde
gold hier alleen voor \LeftR{\exists}. De eigenvariabele~$a$ komt niet
voor in de onderste sequent van die regel; dat is de eindsequent.
Daarom is dit een correcte !!{derivation}.
\end{ex}

\begin{prob}
Geef !!{derivation}s van de volgende sequenten:
\begin{enumerate}
\item $\Sequent (\lforall[x][!A(x)] \land \lforall[y][!B(y)]) \lif
\lforall[z][(!A(z) \land !B(z))]$.
\item $\Sequent (\lexists[x][!A(x)] \lor \lexists[y][!B(y)]) \lif
\lexists[z][(!A(z) \lor !B(z))]$.
\item $\lforall[x][(!A(x) \lif !B)] \Sequent \lexists[y][!A(y)] \lif !B$.
\item $\lforall[x][\lnot !A(x)] \Sequent \lnot\lexists[x][!A(x)]$.
\item $\Sequent \lnot\lexists[x][!A(x)] \lif \lforall[x][\lnot !A(x)]$.
\item $\Sequent \lnot\lexists[x][\lforall[y][((!A(x,y) \lif \lnot
!A(y,y)) \land (\lnot !A(y,y) \lif !A(x,y)))]]$.
\end{enumerate}
\end{prob}

\begin{prob}
Geef !!{derivation}s van de volgende sequenten:
\begin{enumerate}
\item $\Sequent \lnot\lforall[x][!A(x)] \lif \lexists[x][\lnot!A(x)]$.
\item $(\lforall[x][!A(x)] \lif !B) \Sequent \lexists[y][(!A(y) \lif !B)]$.
\item $\Sequent \lexists[x][(!A(x) \lif \lforall[y][!A(y)])]$.
\end{enumerate}
(Bij elk van deze opgaven is de regel $\RightR{\Contraction}$ nodig.)
\end{prob}
\end{document}
